\section{Marco básico de los flujos regularizados continuamente (CNF)}

\subsection{Tres conceptos clave}

El marco matemático del \textbf{Flow Matching} se desarrolla alrededor de tres conceptos centrales, que mantienen estrechas interrelaciones:

\begin{importantbox}{Los tres conceptos clave del Flow Matching}
\begin{enumerate}
    \item \textbf{Mapa de flujo}（Flow Map）$\phi_t(x)$: posición en el instante $t$ a lo largo de la trayectoria que parte desde $x$
    \item \textbf{Campo vectorial}（Vector Field）$v_t(x)$: define la dirección y magnitud de la velocidad en cada punto espacio-temporal dentro de la ODE
    \item \textbf{Trayectoria de densidad de probabilidad}（Probability Density Path）$p_t(x)$: distribución de probabilidad de los puntos de datos en el instante $t$
\end{enumerate}
La relación entre estos tres elementos atraviesa toda la teoría del Flow Matching.
\end{importantbox}

\subsection{Mapa de flujo (Flow Map)}

El mapa de flujo $\phi_t(x_0)$ representa la posición alcanzada tras un tiempo $t$, partiendo desde el punto inicial $x_0$ (procedente de la distribución de referencia $p_0$). Su definición formal es:

$$\phi_t(x_0) = x_0 + \int_0^t v_s(\phi_s(x_0)) \, ds$$

Donde las definiciones de cada símbolo son:
\begin{itemize}
    \item $\phi_t(x_0)$: posición en la trayectoria a partir de $x_0$ en el instante $t$
    \item $v_s(\cdot)$: campo vectorial (campo de velocidades) en el instante $s$
    \item $\phi_0(x_0) = x_0$: condición de frontera; el punto de inicio de la trayectoria cuando $t=0$ es $x_0$ mismo
    \item $\phi_1(x_0) \sim p_1$: condición de frontera; en $t=1$ debe llegar a una muestra de la distribución de datos
\end{itemize}

\begin{knowledgebox}{Relación entre ODE y mapa de flujo}
El mapa de flujo es esencialmente la solución a la siguiente ecuación diferencial ordinaria (ODE):
$$\frac{d\phi_t(x_0)}{dt} = v_t(\phi_t(x_0)), \quad \phi_0(x_0) = x_0$$
Aquí \textbf{no hay aleatoriedad}: una vez fijados el punto de partida $x_0$ y el campo vectorial $v_t$, la trayectoria completa queda determinada. Esto es también lo que significa "Ordinary" en Ordinary Differential Equation, en contraste con los términos estocásticos presentes en las SDE.
\end{knowledgebox}

A veces, para simplificar la notación, utilizamos $x_t$ para denotar $\phi_t(x_0)$, es decir, el punto de la trayectoria a partir de $x_0$ en el instante $t$.

\subsection{Campo vectorial (Vector Field)}

El campo vectorial $v_t(x)$ define la "velocidad" en cada punto espacio-temporal $(t, x)$. Es la derivada temporal del mapa de flujo:

$$v_t(\phi_t(x_0)) = \frac{d\phi_t(x_0)}{dt}$$

El campo vectorial codifica la información del mapeo desde el ruido hacia los datos. Si logramos encontrar un buen campo vectorial que permita mapear muestras extraídas de $p_0$ hacia $p_1$ (distribución de datos), entonces habremos obtenido un modelo generativo.

\subsection{Trayectoria de densidad de probabilidad (Probability Density Path)}

La trayectoria de densidad de probabilidad $p_t(x)$ describe la densidad de probabilidad del punto de datos $x$ en el instante $t$. Satisface las siguientes propiedades:
\begin{itemize}
    \item $p_0 = \mathcal{N}(0, I)$ (distribución de referencia, usualmente una gaussiana estándar)
    \item $p_1 \approx p_{\text{data}}$ (distribución de datos)
    \item En instantes intermedios, $p_t$ es una interpolación entre $p_0$ y $p_1$
\end{itemize}

\begin{knowledgebox}{Ecuación de continuidad (Continuity Equation)}
La relación entre el campo vectorial y la trayectoria de densidad de probabilidad se da mediante la \textbf{ecuación de continuidad}:
$$\frac{\partial p_t(x)}{\partial t} + \nabla \cdot (p_t(x) v_t(x)) = 0$$
Esta es una forma simplificada de la ecuación de Fokker-Planck en ausencia de términos de difusión (es decir, para ODE en lugar de SDE). Establece una conexión matemática exacta entre el campo vectorial y la trayectoria probabilística.
\end{knowledgebox}

\subsection{Relación entre los tres elementos}

En los modelos de difusión, nuestro camino de deducción es:
$$\text{Trayectoria de densidad } p_t \;\longrightarrow\; \text{Forma SDE/ODE} \;\longrightarrow\; \text{Campo vectorial } v_t \;\longrightarrow\; \text{Mapa de flujo } \phi_t$$

Mientras que en el Flow Matching, seguimos un \textbf{camino inverso}:
$$\text{Campo vectorial } v_t \;\longrightarrow\; \text{Mapa de flujo } \phi_t \;\longrightarrow\; \text{Trayectoria de densidad } p_t$$

Es decir, primero definimos el campo vectorial y luego deducimos a partir de él tanto el mapa de flujo como la trayectoria probabilística.

\subsection{Cálculo de la probabilidad logarítmica}

El campo vectorial también puede utilizarse para calcular la probabilidad logarítmica de los puntos de datos. Para $x_1 \sim p_1$, su verosimilitud logarítmica se calcula mediante:

$$\log p_1(x_1) = \log p_0(x_0) - \int_0^1 \nabla \cdot v_t(\phi_t(x_0)) \, dt$$

Donde los términos tienen el siguiente significado:
\begin{itemize}
    \item $\log p_0(x_0)$: probabilidad logarítmica del punto de inicio bajo la distribución de referencia (distribución gaussiana, con forma analítica conocida)
    \item $\nabla \cdot v_t$: divergencia del campo vectorial
    \item Término de integración: acumulación de cambios de volumen a lo largo de la trayectoria
\end{itemize}

Esto implica que si podemos encontrar un campo vectorial que maximice $\log p_1(x_1)$, podremos entrenar un buen modelo generativo.

\subsection{Resumen de la sección}

El marco de los flujos regularizados continuamente (CNF) se compone de tres conceptos centrales: mapa de flujo, campo vectorial y trayectoria de densidad de probabilidad. Estos están interconectados a través de las ODE y la ecuación de continuidad. A diferencia de los modelos de difusión, el CNF define el proceso generativo partiendo directamente del campo vectorial, proporcionando así un marco matemático más directo.


